\section{Data}\label{sec:data}

\subsection{Source and retrieval}

Season-level statistics were retrieved from the \emph{Advanced} tables of Basketball-Reference for every NBA regular season from 1979--80 to 2025--26 \citep{bbref2026}. The tables were downloaded on 4 September 2026 with a documented scraping script (\texttt{code/00\_scrape\_bbref.js}) that records the retrieval date; the raw download (\nRawRows{} player--season--team rows) is released with the paper so that the construction of the analysis sample can be repeated without network access. Each player--season carries PER, age, team, listed position, games and minutes, Win Shares (WS), Box Plus-Minus (BPM) and end-of-season honours (All-Star, All-NBA, MVP voting rank). Players who changed team within a season appear with one aggregate row and one row per team; only the aggregate row was retained, giving \nPlayerSeasons{} player--seasons for \nAllPlayers{} players.

PER is used as published by Basketball-Reference, that is, Hollinger's formula with the site's league- and team-level pace adjustment \citep{hollinger2002,bbrefper}. The earlier version computed PER from box-score totals obtained through \texttt{nba\_api} \citep{nbaapi} with an implementation that omitted the pace adjustment; the reference implementation makes the values comparable with the literature.

\subsection{Sample construction}

The analysis population comprises players whose first NBA season was between 1979--80 and 2010--11. The lower bound marks the introduction of the three-point line and the modern box score; the upper bound makes careers complete or nearly so (only \nActive{} players in the final sample were still active in 2025--26, all with at least \minActiveSeason{} valid seasons). Three further restrictions were applied and are examined in the sensitivity analyses (Section~\ref{sec:sensitivity}). First, a season enters a player's trajectory as a \emph{valid observation} only if the player logged at least 100 minutes; PER computed on a handful of minutes is extremely noisy (in this data the standard deviation of PER is \sdPerLowMp{} for seasons with fewer than 50 minutes against \sdPerHighMp{} for seasons with 500 minutes or more) and would otherwise dominate the reconstruction loss and the distance calculations. Seasons below the threshold are treated as missing, not deleted, so that career length is preserved. Second, players with fewer than 42 career games (half a season) were excluded. Third, at least two valid seasons were required. Table~\ref{tab:flow} gives the resulting flow: of \nDebutWindow{} players who debuted in the window, \nPlayers{} met all criteria, contributing \nObs{} valid PER observations.

\subsection{Representation of a career}

Every retained player $i$ is represented on a calendar axis running from the first to the last valid season, $t=1,\dots,T_i$, with $T_i$ between \minSpan{} and \maxSpan{} seasons. A season inside that span in which the player did not reach 100 minutes (injury, assignment abroad or to the minor league, retirement and return) is a gap: \nGapPlayers{} players (\pctGapPlayers\%) have at least one gap and \nGapSeasons{} of the \nSpanSeasons{} season slots (\pctGapSeasons\%) are gaps. Gaps are encoded by a binary observed-indicator $m_{i,t}$ and are masked in the loss and in the encoder attention; they are not interpolated. This differs from the earlier version, which linearly interpolated interior gaps. The alternative of deleting gaps and treating the remaining seasons as consecutive is evaluated as a sensitivity analysis.

PER values were standardised with the mean (\muPer) and standard deviation (\sdPer) of all valid observations in the development set (defined below) before entering the autoencoder; reconstruction errors are reported after inverting the transformation, in PER units. No other transformation was applied: PER is already normalised so that the league average equals 15 in every season, and its distribution across valid observations is close to symmetric (skewness \skewPer).

\subsection{External variables}

Variables used only for validation, never for clustering, were: the modal listed position over valid seasons (PG, SG, SF, PF, C); draft status, obtained by matching names and debut years to the NBA draft history retrieved through \texttt{nba\_api} (lottery pick 1--14, other first round, second round, undrafted; \nDrafted{} of \nPlayers{} players matched to a draft record, the remainder classified as undrafted); whether the player was ever selected to an All-Star game or an All-NBA team; total career WS and mean BPM over valid seasons. The BPM trajectories of the same players and seasons were also used to repeat the whole analysis with a different performance metric.
